\lecture{Fondamenti di Impianti Industriali}{23~-~01~-~2026}
%
\section{Definizione e compentenze}
Per affrontare correttamente la materia, è necessario sviluppare una visione d'insieme che vada oltre la semplice definizione. L'obiettivo è acquisire la capacità di progettare, gestire e ottimizzare un sito industriale. Nello specifico, bisogna essere in grado di:
\begin{itemize}
    \item Comprendere e rappresentare il processo tecnologico e l'impianto nel suo insieme.
    \item Dimensionare non solo gli impianti di processo, ma anche quelli tecnologici e di servizio ausiliari, fondamentali per il funzionamento dello stabilimento.
    \item Costruire il flusso di processo, integrando logistica interna ed esterna, produzione e stoccaggio.
    \item Elaborare uno Studio di Fattibilità e stimare il Budget economico necessario alla realizzazione.
    \item Gestire i rischi per la salute e la sicurezza dei lavoratori.
    \item Utilizzare tecniche moderne come la virtualizzazione tramite Digital Twin per simulare l'impianto.
\end{itemize}
%
\subsection{Manufacturing e sistema di produzione}
Il settore Manifatturiero (\emph{Manufacturing}) è il cuore dell'attività industriale. Esso rappresenta l'attività umana volta a trasformare le materie prime in:
\begin{enumerate}
    \item Prodotti finiti: destinati al consumatore finale.
    \item Semilavorati: destinati a ulteriori lavorazioni manifatturiere.
\end{enumerate}
Il Sistema di Produzione funziona secondo un modello di input-output ben definito:
\begin{itemize}
    \item \emphr{Input}: Materia prima/semilavorati, Energia, Personale, Fattori di produzione.
    \item \emphr{Processo}: Una ``scatola nera'' centrale dove avviene la trasformazione.
    \item \emphr{Output}: Prodotti finiti/semilavorati.
\end{itemize}
La fase cruciale di trasformazione (sia essa fisica o chimica) che permette di ottenere nuovi beni in uscita è definita \emph{Processo Tecnologico}.
%
\subsection{Impianto Industriale}
L'impianto industriale è un complesso sistema tecnologico che include macchinari, attrezzature e processi organizzati per la produzione su larga scala. Esso si inserisce sempre all'interno di un'impresa con obiettivi economici.
In ambito tecnico e normativo, l'impianto può essere indicato con vari sinonimi: Unità produttiva, Stabilimento, Fabbrica, Officina, Opificio, Laboratorio, Plant \dots
%
\paragraph{Composizione strutturale dello stabilimento produttivo} Uno stabilimento non è un monolite, ma si divide in due macro-categorie che operano in sinergia:
\begin{enumerate}
    \item Impianto Tecnologico o Produttivo: Dove avviene la trasformazione diretta.
    \item Impianti Ausiliari o di Servizio: Necessari al supporto e al funzionamento del reparto produttivo.
\end{enumerate}
%
\paragraph{Criteri di Sviluppo} La realizzazione di un impianto deve seguire criteri etici ed economici precisi:
\begin{itemize}
    \item Convenienza per la collettività e per l'azienda.
    \item Rispetto dell'ambiente e annullamento dell'inquinamento.
    \item Salvaguardia della salute e sicurezza dei lavoratori.
    \item Risparmio energetico e uso razionale delle risorse.
\end{itemize}
%
\section{Supply Chain}
L'impianto industriale non è un'isola, ma il punto centrale della Supply Chain (catena di approvvigionamento). La Supply Chain è il processo end-to-end che va dalle materie prime alla consegna al consumatore.
%
\paragraph{Integrazione strategica} L'impianto funge da collegamento cruciale tra fornitori e distributori. Per il successo della catena, è necessaria una collaborazione sinergica su diversi fronti:
\begin{itemize}
    \item Con i Fornitori: Scambio tempestivo di informazioni per evitare ritardi nell'arrivo delle materie prime.
    \item Gestione dell'Inventario: Coordinamento degli stock per evitare sovrapproduzione o carenze (out-of-stock), utilizzando sistemi di tracciabilità.
    \item Adattabilità: Capacità di regolare la produzione basandosi su previsioni e analisi dei dati per rispondere alle fluttuazioni della domanda di mercato.
    \item Comunicazione: Scambio regolare di dati con tutti i partner chiave.
\end{itemize}
%
\section{Classificazione degli Impianti Industriali}
Gli impianti industriali vengono classificati secondo due criteri principali: la natura del processo produttivo e le dimensioni economiche.

%
\paragraph{Classificazione per Natura del Processo}
\begin{itemize}
    \item Manifatturiero: Trasformazione in prodotti finiti (es. auto, tessile).
    \item Estrattivo: Miniere e cave.
    \item Energetico: Produzione di energia (termica, solare, eolica, ecc.).
    \item Chimico e Petrolchimico: Processi chimici complessi e raffinazione di combustibili.
    \item Alimentare: Trattamento e confezionamento cibi.
    \item Farmaceutico: Produzione di medicinali.
    \item Trattamento delle acque: Depurazione reflui.
    \item Altri settori: Edile, Elettronico, Automobilistico, Tessile \dots
\end{itemize}

%
\paragraph{Classificazione dimensionale (per criteri EU)}
La classificazione delle imprese si basa sul numero di dipendenti e sul fatturato annuo (o totale di bilancio), criteri stabiliti dall'UE. Vedi tabella \ref{tab:class_imp} a riguardo.
\begin{table}
    \centering
    \begin{tabular}{lrr}%
        \toprule
        \textbf{Tipologia} & \textbf{Dipententi} & \textbf{Fatturato o totale bilancio} \\
        \midrule
        Microimpresa       & $ <10 $ \%          & $ \leq 2 $ Milioni di Euro           \\
        Piccola Impresa    & $ <50 $ \%          & $ \leq 10 $ Milioni di Euro          \\
        Grande Impresa     & $ >250 $ \%         & $ \leq 50 $ Milioni di Euro          \\
        Media Impresa      & $ <250 $ \%         & $ >50 $ Milioni di Euro              \\
        \bottomrule
    \end{tabular}
    \caption{Classificazione delle Imprese in UE.}\label{tab:class_imp}
\end{table}
